%%
%% The abstract is a short summary of the work to be presented in the
%% article.
\begin{abstract}
    EDA scripting with tool-specific, often undocumented APIs remains a long-tail 
    bottleneck that existing LLM-based tooling fails to address. This paper presents 
    \textbf{ZhuLong}, an execution-grounded LLM coding agent for PyAether, SKILL, and Tcl 
    that combines API retrieval, documentation inspection, and sandbox execution via 
    unified MCP tools into a closed generate--execute--refine loop. We further provide 
    a \textbf{theoretical grounding} for why runtime observation is necessary: EDA correctness 
    is a \emph{situated} property that cannot be guaranteed from static knowledge, 
    and we formalize this ``grounding gap'' with a Bayesian argument showing it is 
    bounded by observation fidelity---not model scale. The Harness supplies precisely 
    that observation, by reading the environment's documented API facts and by reading 
    back what executing code does to the live design.
    % 使用工具特定且常无文档的EDA API进行脚本编写，仍是现有LLM工具无法解决的长尾瓶颈。
    % 本文提出ZhuLong，一个面向PyAether和SKILL的执行驱动LLM代码智能体，通过统一MCP工具将
    % API检索、文档查阅和沙盒执行组合为"生成-执行-修正"闭环。我们进一步给出执行必要性的
    % 理论论证：EDA正确性是情境化属性，无法从静态知识保证；我们用贝叶斯论证形式化这一
    % "落地鸿沟"，说明它受观测保真度而非模型规模约束，沙盒执行正是弥合该鸿沟的观测通道。

    We evaluate ZhuLong on \textbf{EDA-Eval-PyAether}, a benchmark of 158 real-world 
    tasks with assertion-based execution, where the complete system achieves 
    \textbf{84.8\%} Pass@1 in the commercial Empyrean Aether environment, substantially 
    outperforming a pure LLM baseline (\textbf{10.5\%}). Ablation studies show the 
    execution loop is decisive: adding sandbox execution to retrieval yields +[TBD] pp, 
    while removing either execution ($-[TBD]$ pp) or retrieval ($-3.8$ pp) reduces 
    accuracy---consistent with the theoretical prediction.
    % 我们在EDA-Eval-PyAether（158个真实任务、断言式执行验证）上评估ZhuLong，完整系统在
    % 商业Empyrean Aether环境中达到84.8\%的Pass@1，显著优于纯LLM基线（10.5\%）。消融表明
    % 执行闭环是决定性的：在检索之上加入沙盒执行带来+[TBD]个百分点，而移除执行（-[TBD]pp）
    % 或检索（-3.8pp）都会使准确率坍塌——与理论预测一致。
All quantitative results in this draft are provisional, obtained on the PyAether-only 
    slice; final numbers will be reported on the unified \textbf{EDA-Scripting-Bench} 
    (PyAether + SKILL + Tcl).
    % 本文所有定量结果为暂定值，均来自PyAether切片；最终结果将在统一评测集EDA-Scripting-Bench（PyAether + SKILL + Tcl）上报告。

\end{abstract}

\keywords{LLM Agents for EDA, EDA Scripting, Code Generation, Sandbox Execution, Execution Grounding}